\section{Méthodes et protocole de validation}

\paragraph{Une régression par propriété.} Chaque grandeur a son modèle, entraîné sur
les lignes où elle est mesurée. Pour le Charpy, la température d'essai est une
\emph{entrée} : la même soudure absorbe \SI{39}{\joule} à \SI{-60}{\celsius} et
\SI{114}{\joule} à \SI{0}{\celsius}. Ce n'est pas un levier : elle est imposée par
l'usage de la pièce et reste fixée dans les recommandations.

\paragraph{Protocole d'évaluation.}\label{sec:f-protocole} Tous les modèles suivent le
même protocole, pour être comparables \cite{hastie2009,pedregosa2011} :
\begin{itemize}
  \item \textbf{validation croisée groupée par article} (\code{GroupKFold}, 5 plis) :
    toutes les lignes d'un article sont dans le même pli. Un découpage au hasard
    mettrait des soudures quasi identiques en entraînement et en test, et serait trop
    optimiste. Le score mesure la capacité à prédire une \emph{étude inconnue} ;
  \item \textbf{validation croisée imbriquée} \cite{varma2006} : dans chaque pli, une
    boucle intérieure (elle aussi par article) choisit les hyperparamètres sur la seule
    partie entraînement ; le pli de test ne sert jamais au réglage ;
  \item \textbf{pré-traitement dans un \code{Pipeline}} : imputation, encodage et
    standardisation sont réappris dans chaque pli, sans fuite d'information.
\end{itemize}

\paragraph{Mesures d'erreur.} Chaque ligne est prédite une fois, par un modèle qui ne l'a
pas vue. Sur les $n$ couples (valeur réelle $y_i$, prédiction $\hat y_i$) :
\[
  \mathrm{RMSE} = \sqrt{\tfrac{1}{n}\textstyle\sum_i (y_i - \hat y_i)^2}, \qquad
  \mathrm{MAE} = \tfrac{1}{n}\textstyle\sum_i \lvert y_i - \hat y_i \rvert, \qquad
  R^2 = 1 - \dfrac{\sum_i (y_i - \hat y_i)^2}{\sum_i (y_i - \bar y)^2}.
\]
La RMSE, dans l'unité de la cible, pèse fortement les grosses erreurs, les plus graves
ici (rater une soudure fragile) ; c'est aussi ce que minimise le réglage. La MAE donne
l'erreur typique. Le $R^2$ compare le modèle à une prédiction constante égale à la
moyenne : 0 signifie « pas mieux que la moyenne », un $R^2$ négatif « pire ». Il permet
de comparer des cibles d'unités différentes. Les mesures sont calculées sur toutes les
prédictions réunies, car un $R^2$ par pli serait instable ; l'écart-type de la RMSE entre
plis mesure la stabilité.

\paragraph{Modèles comparés.}
\begin{itemize}
  \item \textbf{Référence} : la moyenne de l'entraînement, le minimum à battre.
  \item \textbf{Régression linéaire}, puis \textbf{régression sur composantes principales
    (PCR) et PLS} \cite{wold2001}, qui résument les entrées corrélées en quelques axes
    (nombre d'axes réglé par validation croisée).
  \item \textbf{Arbre de régression} \cite{breiman1984}, \textbf{bagging}
    \cite{breiman1996} et \textbf{forêt aléatoire} \cite{breiman2001} : un arbre gère
    seuils et interactions et n'extrapole pas, mais a une forte variance ; le bagging
    moyenne 300 arbres entraînés sur des tirages des lignes, et la forêt tire en plus
    une partie des variables à chaque découpage, ce qui décorrèle les arbres. Réglages :
    profondeur et taille des feuilles, part des variables tirées.
  \item \textbf{Semi-supervisé} \cite{zhu2009} : pour chaque cible, 47 à 58\,\% des lignes
    ont des entrées mais pas de mesure. L'\emph{auto-apprentissage} réentraîne la forêt
    avec ses propres prédictions sur ces lignes, en ne gardant que celles où ses arbres
    s'accordent. La co-régression \emph{COREG} \cite{zhou2005}, issue du co-apprentissage
    \cite{blum1998}, entraîne deux régresseurs des plus proches voisins sur deux vues
    (composition ; procédé et traitement) qui s'étiquettent mutuellement ces lignes. Ces
    deux méthodes sont les plus simples à adosser à nos modèles, et COREG est conçue
    pour la régression.
  \item \textbf{Verdict « bonne / mauvaise »} : les prédictions hors pli de la forêt sont
    comparées aux seuils de la classe ISO 2560 E~42~3 \cite{iso2560} : limite
    d'élasticité $\geq$ \SI{420}{\mega\pascal}, allongement $\geq$ \SI{20}{\percent},
    Charpy $\geq$ \SI{47}{\joule} à \SI{-30}{\celsius}. La soudure est bonne si elle
    passe les trois.
\end{itemize}
